\documentclass{article}
\usepackage[margin=0.8in]{geometry}
\usepackage{amsmath,amssymb,amsthm,mathtools}
\usepackage{mathrsfs,bm,bbm}
\usepackage{graphicx,booktabs,array,multirow,tabularx,longtable}
\usepackage{enumitem}
\usepackage{xcolor}
\usepackage{algorithm,algpseudocode}
\usepackage{subcaption,tikz}
\usepackage{siunitx,cancel,accents}
\usepackage{microtype}
\usepackage[numbers,sort&compress]{natbib}
\usepackage[colorlinks=true,linkcolor=blue,citecolor=blue,urlcolor=blue]{hyperref}
\usepackage{cleveref}
\setlength{\emergencystretch}{2em}
\newtheorem{assumption}{Assumption}
\newtheorem{definition}{Definition}
\newtheorem{theorem}{Theorem}
\newtheorem{lemma}{Lemma}
\newtheorem{proposition}{Proposition}
\newtheorem{openendedquestion}{Open-ended Question}
\newtheorem{conjecture}{Conjecture}
\newcommand{\coreref}[1]{\ref{#1}}

\begin{document}

\sloppy

\section*{stat\_\allowbreak{}plm\_\allowbreak{}two\_\allowbreak{}nuisance\_\allowbreak{}minimax}

\noindent\textit{Specialization:} v1\par

\paragraph{TL;DR.} When \(\nu_d=\delta_{x_0}\), both nuisance functions are constants almost surely and \(\beta(P)=\operatorname{Cov}_P(T,Y)/\operatorname{Var}_P(T)\). An explicit clipped covariance estimator and a bounded non-atomic Le Cam submodel prove the matching minimax order \(m_{\rm PL}^{0}(n,\bar\delta_\mu,\bar\delta_\pi)\asymp n^{-1/2}\) for every budget pair. Thus product-scale behavior cannot hold uniformly over the frozen design class without a quantitative covariate-richness condition. Appendix A.1 of \cite{gu2026optimally} uses the outcome-error subscript \(\varepsilon_Y\); this note makes no source-typo claim and does not tighten its displayed lower bound.

\section{Project justification}

\paragraph{Gap.} Gu, Yin, Cai, and Fan's two-nuisance displays do not determine which covariate designs can sustain a product term, while the frozen core imposes no quantitative richness condition on \(\nu_d\).

\paragraph{Niche.} The note isolates an admissible boundary design on which both nuisance functions collapse to constants and determines the matching root-\(n\) minimax order in the same fixed-class experiment.

\paragraph{Fill.} A full-model covariance upper bound and a bounded non-atomic two-point lower bound prove \(m_{\rm PL}^{0}\asymp n^{-1/2}\) when \(\nu_d=\delta_{x_0}\). This is a covariate-richness obstruction and does not tighten Gu, Yin, Cai, and Fan's displayed lower bound.

\section{Related work}

Robinson \cite{robinson1988root} gives the classical root-\(n\) partially linear benchmark, and Chernozhukov et al. \cite{chernozhukov2018dml} give orthogonal-score inference for the coefficient. Gu \cite{gu2025open} identifies variance propagation as an unresolved structure-agnostic issue. Gu, Yin, Cai, and Fan \cite{gu2026optimally}, Appendix A.1, study the same fixed-class, two-nuisance, \(2n\)-sample experiment and display nonmatching upper and lower expressions. The present theorem specializes to an allowed degenerate reference law, supplies a covariate-richness obstruction, and does not tighten their displayed lower bound. Jin, Mackey, and Syrgkanis \cite{jin2025hardnormal} study supplied nuisance estimates and worst-case quantile loss under binary or conditionally Gaussian treatment, while Jin and Syrgkanis \cite{jin2025sharp} develop lower-bound methods for other functionals. Balakrishnan, Kennedy, and Wasserman \cite{balakrishnan2023fundamental} provide the broader structure-agnostic lower-bound perspective.

\section{Setup}

Target (estimand / structural parameter): \(\beta(P)\).
Primary functional / estimator / diagnostic: \par\noindent \makebox[\linewidth][c]{\resizebox{0.98\linewidth}{!}{\(\displaystyle m_{\rm PL}^{0}(n,\bar\delta_{\mu,n},\bar\delta_{\pi,n})=\sup_{F_\mu\in\mathcal H_n(\bar\delta_{\mu,n})}\sup_{F_\pi\in\mathcal H_n(\bar\delta_{\pi,n})}\inf_{\widehat\beta}\sup_{P\in\mathcal P_{\rm PL}^{0}(F_\mu,F_\pi)}\int|\widehat\beta-\beta(P)|\,dP^{\otimes2n}\)}}\par\noindent.
\begin{itemize}
  \item \(n\) --- \texttt{integer}; \(\mathbb N\); \texttt{index}
  \par\smallskip\noindent\emph{Definition.} Half-sample-size index; the experiment observes \(2n\) units.
  \item \(d\) --- \texttt{integer}; \(\mathbb N\); \texttt{index}
  \par\smallskip\noindent\emph{Definition.} \texttt{Covariate dimension.}
  \item \(\nu_d\) --- \texttt{probability law}; \(\mathcal P([0,1]^d)\); \texttt{reference law}
  \par\smallskip\noindent\emph{Definition.} \texttt{Reference covariate law.}
  \item \(\mathcal T_3\) --- \texttt{function class}; \(L^2(\nu_d)\); \texttt{ambient class}
  \par\smallskip\noindent\emph{Definition.} Ambient class \(\mathcal T_3=\{f\in L^2(\nu_d):\lVert f\rVert_\infty\le3\}\).
  \item \(F\) --- \texttt{function class}; \(2^{\mathcal T_3}\); \texttt{truth class}
  \par\smallskip\noindent\emph{Definition.} \texttt{Generic truth class.}
  \item \(G\) --- \texttt{function class}; \(2^{\mathcal T_3}\); \texttt{learner witness}
  \par\smallskip\noindent\emph{Definition.} Generic learner-class witness for \(F\).
  \item \(\partial G\) --- \texttt{function class}; \(\{g-g':g,g'\in G\}\); \texttt{complexity class}
  \par\smallskip\noindent\emph{Definition.} Difference class of \(G\).
  \item \(\mathcal R_{n,\nu_d}\) --- \texttt{functional}; \(\mathbb R_+\times2^{L^2(\nu_d)}\to\mathbb R_+\); \texttt{complexity functional}
  \par\smallskip\noindent\emph{Definition.} \texttt{Localized Rademacher-\allowbreak{}complexity functional.}
  \item \(\delta^{\rm appr}\) --- \texttt{nonnegative scalar}; \([0,1]\); \texttt{budget}
  \par\smallskip\noindent\emph{Definition.} \texttt{Approximation budget.}
  \item \(\delta^{\rm stoc}\) --- \texttt{nonnegative scalar}; \([0,1]\); \texttt{budget}
  \par\smallskip\noindent\emph{Definition.} \texttt{Stochastic budget.}
  \item \(\bar\delta\) --- \texttt{budget pair}; \([0,1]^2\); \texttt{budget}
  \par\smallskip\noindent\emph{Definition.} Pair \(\bar\delta=(\delta^{\rm appr},\delta^{\rm stoc})\).
  \item \(\delta_\mu^{\rm appr}\) --- \texttt{nonnegative scalar}; \([0,1]\); \texttt{budget}
  \par\smallskip\noindent\emph{Definition.} \texttt{Outcome approximation component.}
  \item \(\delta_\mu^{\rm stoc}\) --- \texttt{nonnegative scalar}; \([0,1]\); \texttt{budget}
  \par\smallskip\noindent\emph{Definition.} \texttt{Outcome stochastic component.}
  \item \(\delta_\pi^{\rm appr}\) --- \texttt{nonnegative scalar}; \([0,1]\); \texttt{budget}
  \par\smallskip\noindent\emph{Definition.} \texttt{Treatment-\allowbreak{}regression approximation component.}
  \item \(\delta_\pi^{\rm stoc}\) --- \texttt{nonnegative scalar}; \([0,1]\); \texttt{budget}
  \par\smallskip\noindent\emph{Definition.} \texttt{Treatment-\allowbreak{}regression stochastic component.}
  \item \(\mathcal H_n(\bar\delta)\) --- \texttt{class of function classes}; \(2^{2^{\mathcal T_3}}\); \texttt{admissible truth-\allowbreak{}class collection}
  \par\smallskip\noindent\emph{Definition.} Truth classes with a uniform learner-class witness at budget \(\bar\delta\).
  \item \(F_\mu\) --- \texttt{function class}; \(2^{\mathcal T_3}\); \texttt{truth class}
  \par\smallskip\noindent\emph{Definition.} \texttt{Fixed structural-\allowbreak{}outcome truth class.}
  \item \(F_\pi\) --- \texttt{function class}; \(2^{\mathcal T_3}\); \texttt{truth class}
  \par\smallskip\noindent\emph{Definition.} \texttt{Fixed treatment-\allowbreak{}regression truth class.}
  \item \(G_\mu\) --- \texttt{function class}; \(2^{\mathcal T_3}\); \texttt{learner witness}
  \par\smallskip\noindent\emph{Definition.} Learner witness for \(F_\mu\).
  \item \(G_\pi\) --- \texttt{function class}; \(2^{\mathcal T_3}\); \texttt{learner witness}
  \par\smallskip\noindent\emph{Definition.} Learner witness for \(F_\pi\).
  \item \(X\) --- \texttt{random vector}; \([0,1]^d\); \texttt{covariate}
  \par\smallskip\noindent\emph{Definition.} \texttt{Observed covariates.}
  \item \(T\) --- \texttt{random variable}; \([-3,3]\); \texttt{treatment}
  \par\smallskip\noindent\emph{Definition.} \texttt{Observed treatment.}
  \item \(Y\) --- \texttt{random variable}; \([-3,3]\); \texttt{outcome}
  \par\smallskip\noindent\emph{Definition.} \texttt{Observed outcome.}
  \item \(u\) --- \texttt{random variable}; \(\mathbb R\); \texttt{error}
  \par\smallskip\noindent\emph{Definition.} Treatment innovation in \(T=\pi(X)+u\).
  \item \(\varepsilon\) --- \texttt{random variable}; \(\mathbb R\); \texttt{error}
  \par\smallskip\noindent\emph{Definition.} Mean-zero outcome error in \(Y=\beta T+\mu(X)+\varepsilon\).
  \item \(\mu\) --- \texttt{measurable function}; \(F_\mu\); \texttt{nuisance}
  \par\smallskip\noindent\emph{Definition.} \texttt{Structural-\allowbreak{}outcome nuisance.}
  \item \(\pi\) --- \texttt{measurable function}; \(F_\pi\); \texttt{nuisance}
  \par\smallskip\noindent\emph{Definition.} \texttt{Treatment regression.}
  \item \(\beta\) --- \texttt{real scalar}; \([-3,3]\); \texttt{target parameter}
  \par\smallskip\noindent\emph{Definition.} \texttt{Partially linear coefficient.}
  \item \(v\) --- \texttt{real scalar}; \([1/3,3]\); \texttt{variance}
  \par\smallskip\noindent\emph{Definition.} \texttt{Conditional treatment-\allowbreak{}innovation variance.}
  \item \(P_{\beta,v,\mu,\pi}^0\) --- \texttt{probability law}; \(\mathcal P([0,1]^d\times[-3,3]^2)\); \texttt{law}
  \par\smallskip\noindent\emph{Definition.} \texttt{Mean-\allowbreak{}zero partially linear observed-\allowbreak{}data law.}
  \item \(P\) --- \texttt{probability law}; \(\mathcal P([0,1]^d\times[-3,3]^2)\); \texttt{law}
  \par\smallskip\noindent\emph{Definition.} \texttt{Generic law in the mean-\allowbreak{}zero fixed-\allowbreak{}class model.}
  \item \(\mathcal P_{\rm PL}^{0}(F_\mu,F_\pi)\) --- \texttt{statistical model}; \(2^{\mathcal P([0,1]^d\times[-3,3]^2)}\); \texttt{experiment}
  \par\smallskip\noindent\emph{Definition.} \texttt{Mean-\allowbreak{}zero full two-\allowbreak{}nuisance partially linear model for a fixed class pair.}
  \item \(\mathcal P_{\rm PL}^{0,\rm na}(F_\mu,F_\pi)\) --- \texttt{statistical model}; \(2^{\mathcal P([0,1]^d\times[-3,3]^2)}\); \texttt{hard subclass}
  \par\smallskip\noindent\emph{Definition.} Bounded non-atomic-treatment subclass of \(\mathcal P_{\rm PL}^{0}(F_\mu,F_\pi)\).
  \item \(\beta(P)\) --- \texttt{real scalar}; \([-3,3]\); \texttt{target}
  \par\smallskip\noindent\emph{Definition.} Coefficient identified by a law \(P\) in the mean-zero model.
  \item \(\widehat\beta\) --- \texttt{estimator}; \(([0,1]^d\times[-3,3]^2)^{2n}\to\mathbb R\); \texttt{decision rule}
  \par\smallskip\noindent\emph{Definition.} \texttt{Measurable estimator tailored to the selected fixed truth classes and budgets.}
  \item \(\widehat\beta_n\) --- \texttt{estimator}; \(([0,1]^d\times[-3,3]^2)^{2n}\to\mathbb R\); \texttt{decision rule}
  \par\smallskip\noindent\emph{Definition.} An \(n\)-indexed estimator tailored to a selected hard class pair.
  \item \(o_{1:2n}\) --- \texttt{sample point}; \(([0,1]^d\times[-3,3]^2)^{2n}\); \texttt{integration variable}
  \par\smallskip\noindent\emph{Definition.} Realized \(2n\)-sample point.
  \item \(\bar\delta_\mu\) --- \texttt{budget pair}; \([0,1]^2\); \texttt{budget}
  \par\smallskip\noindent\emph{Definition.} \texttt{Generic outcome budget pair.}
  \item \(\bar\delta_\pi\) --- \texttt{budget pair}; \([0,1]^2\); \texttt{budget}
  \par\smallskip\noindent\emph{Definition.} \texttt{Generic treatment-\allowbreak{}regression budget pair.}
  \item \(s_{\mu,n}\) --- \texttt{nonnegative scalar}; \([0,1]\); \texttt{budget}
  \par\smallskip\noindent\emph{Definition.} \texttt{Outcome stochastic budget.}
  \item \(s_{\pi,n}\) --- \texttt{nonnegative scalar}; \([0,1]\); \texttt{budget}
  \par\smallskip\noindent\emph{Definition.} \texttt{Treatment-\allowbreak{}regression stochastic budget.}
  \item \(\bar\delta_{\mu,n}\) --- \texttt{budget pair}; \([0,1]^2\); \texttt{budget}
  \par\smallskip\noindent\emph{Definition.} Outcome budget \(\bar\delta_{\mu,n}=(0,s_{\mu,n})\).
  \item \(\bar\delta_{\pi,n}\) --- \texttt{budget pair}; \([0,1]^2\); \texttt{budget}
  \par\smallskip\noindent\emph{Definition.} Treatment budget \(\bar\delta_{\pi,n}=(0,s_{\pi,n})\).
  \item \(\mathcal R_n^{\rm asym}\) --- \texttt{set of budget sequences}; \(\bigl([0,1]^2\times[0,1]^2\bigr)^{\mathbb N}\); \texttt{regime}
  \par\smallskip\noindent\emph{Definition.} \texttt{Asymmetric pure-\allowbreak{}stochastic budget regime.}
  \item \(m_{\rm PL}^{0}(n,\bar\delta_\mu,\bar\delta_\pi)\) --- \texttt{nonnegative scalar}; \(\mathbb R_+\); \texttt{risk functional}
  \par\smallskip\noindent\emph{Definition.} \texttt{Fixed-\allowbreak{}class minimax expected absolute risk of the mean-\allowbreak{}zero full two-\allowbreak{}nuisance partially linear model.}
  \item \(\mathsf U_n\) --- \texttt{nonnegative scalar}; \(\mathbb R_+\); \texttt{benchmark}
  \par\smallskip\noindent\emph{Definition.} \texttt{Theorem A.1 displayed upper expression in the proposed regime.}
  \item \(\mathsf L_n\) --- \texttt{nonnegative scalar}; \(\mathbb R_+\); \texttt{benchmark}
  \par\smallskip\noindent\emph{Definition.} \texttt{Theorem A.1 displayed lower expression in the proposed regime.}
  \item \(F_{\mu,n}\) --- \texttt{function class}; \(\mathcal H_n(0,s_{\mu,n})\); \texttt{hard truth class}
  \par\smallskip\noindent\emph{Definition.} Fixed outcome truth class in the \(n\)-th hard witness.
  \item \(F_{\pi,n}\) --- \texttt{function class}; \(\mathcal H_n(0,s_{\pi,n})\); \texttt{hard truth class}
  \par\smallskip\noindent\emph{Definition.} Fixed treatment truth class in the \(n\)-th hard witness.
\end{itemize}

\section{Assumptions}

\begin{assumption}[ass:uniform-approximation]\label{ass:uniform-approximation}
A pair \(F,G\subseteq\mathcal T_3\) obeys \(\sup_{f\in F}\inf_{g\in G}\lVert f-g\rVert_{L^2(\nu_d)}\le\delta^{\rm appr}\).
\par\smallskip\noindent\emph{Standard} (Gu--Yin--Cai--Fan uniform approximation condition, \cite{gu2026optimally}).
\end{assumption}

\begin{assumption}[ass:localized-critical-radius]\label{ass:localized-critical-radius}
A learner class \(G\subseteq\mathcal T_3\) obeys \par\noindent \makebox[\linewidth][c]{\resizebox{0.98\linewidth}{!}{\(\displaystyle \sup_{\delta>0:\,\delta\ge\delta^{\rm stoc}}\mathcal R_{n,\nu_d}(\delta;\partial G)/(3\delta)\le\delta^{\rm stoc}\)}}\par\noindent.
\par\smallskip\noindent\emph{Standard} (Gu--Yin--Cai--Fan localized-critical-radius condition, \cite{gu2026optimally}).
\end{assumption}

\begin{assumption}[ass:two-n-iid]\label{ass:two-n-iid}
The \(2n\) observations are independent and identically distributed under \(P\).
\par\smallskip\noindent\emph{Standard} (i.i.d. sampling, \cite{gu2026optimally}).
\end{assumption}

\section{Definitions}

\begin{definition}[def:admissible-truth-classes]\label{def:admissible-truth-classes}
\par\smallskip\noindent\emph{Defined object.} \texttt{Admissible truth classes}
\par\smallskip\noindent\emph{Construction.} \par\noindent \makebox[\linewidth][c]{\resizebox{0.98\linewidth}{!}{\(\displaystyle \mathcal H_n(\bar\delta)=\{F\subseteq\mathcal T_3:\exists G\subseteq\mathcal T_3,\ \sup_{f\in F}\inf_{g\in G}\lVert f-g\rVert_{L^2(\nu_d)}\le\delta^{\rm appr},\ \sup_{\delta>0:\,\delta\ge\delta^{\rm stoc}}\mathcal R_{n,\nu_d}(\delta;\partial G)/(3\delta)\le\delta^{\rm stoc}\}\)}}\par\noindent, where \(\bar\delta=(\delta^{\rm appr},\delta^{\rm stoc})\).
\par\smallskip\noindent\emph{Member properties.} \texttt{ass:\allowbreak{}uniform-\allowbreak{}approximation}; \texttt{ass:\allowbreak{}localized-\allowbreak{}critical-\allowbreak{}radius}.
\end{definition}

\begin{definition}[def:mean-zero-plm-family]\label{def:mean-zero-plm-family}
\par\smallskip\noindent\emph{Defined object.} \texttt{Mean-\allowbreak{}zero full two-\allowbreak{}nuisance partially linear family}
\par\smallskip\noindent\emph{Construction.} \par\noindent \makebox[\linewidth][c]{\resizebox{0.98\linewidth}{!}{\(\displaystyle \mathcal P_{\rm PL}^{0}(F_\mu,F_\pi)=\{P_{\beta,v,\mu,\pi}^{0}:X\sim\nu_d,\ \mu\in F_\mu,\ \pi\in F_\pi,\ T=\pi(X)+u,\ \mathbb E[u\mid X]=0,\ \mathbb E[u^2\mid X]\equiv v\ge1/3,\ Y=\beta T+\mu(X)+\varepsilon,\ \mathbb E[\varepsilon\mid X,T]=0,\ |\beta|\vee|v|\le3,\ |T|\vee|Y|\le3\}\)}}\par\noindent. This is the mean-zero full two-nuisance PLM consistent with equations (1.1) and (3.2) of \cite{gu2026optimally}.
\par\smallskip\noindent\emph{Inputs.} \(F_\mu\); \(F_\pi\); \(\nu_d\); \(X\); \(T\); \(Y\); \(u\); \(\varepsilon\); \(\mu\); \(\pi\); \(\beta\); \(v\).
\end{definition}

\begin{definition}[def:non-atomic-hard-subclass]\label{def:non-atomic-hard-subclass}
\par\smallskip\noindent\emph{Defined object.} \texttt{Bounded non-\allowbreak{}atomic treatment subclass}
\par\smallskip\noindent\emph{Construction.} \par\noindent \makebox[\linewidth][c]{\resizebox{0.98\linewidth}{!}{\(\displaystyle \mathcal P_{\rm PL}^{0,\rm na}(F_\mu,F_\pi)=\{P\in\mathcal P_{\rm PL}^{0}(F_\mu,F_\pi):P(T=t)=0\text{ for every }t\in[-3,3]\}\)}}\par\noindent.
\par\smallskip\noindent\emph{Inputs.} \(\mathcal P_{\rm PL}^{0}(F_\mu,F_\pi)\); \(F_\mu\); \(F_\pi\); \(T\).
\end{definition}

\begin{definition}[def:mean-zero-fixed-class-risk]\label{def:mean-zero-fixed-class-risk}
\par\smallskip\noindent\emph{Defined object.} \texttt{Mean-\allowbreak{}zero fixed-\allowbreak{}class minimax risk}
\par\smallskip\noindent\emph{Construction.} \par\noindent \makebox[\linewidth][c]{\resizebox{0.98\linewidth}{!}{\(\displaystyle m_{\rm PL}^{0}(n,\bar\delta_\mu,\bar\delta_\pi)=\sup_{F_\mu\in\mathcal H_n(\bar\delta_\mu)}\sup_{F_\pi\in\mathcal H_n(\bar\delta_\pi)}\inf_{\widehat\beta}\sup_{P\in\mathcal P_{\rm PL}^{0}(F_\mu,F_\pi)}\int|\widehat\beta(o_{1:2n})-\beta(P)|\,dP^{\otimes2n}(o_{1:2n})\)}}\par\noindent. The infimum occurs after selection of the fixed class pair, so the estimator may know the truth classes and budgets.
\par\smallskip\noindent\emph{Inputs.} \(n\); \(\bar\delta_\mu\); \(\bar\delta_\pi\); \(\mathcal H_n(\bar\delta)\); \(F_\mu\); \(F_\pi\); \(\widehat\beta\); \(\mathcal P_{\rm PL}^{0}(F_\mu,F_\pi)\); \(\beta(P)\); \(o_{1:2n}\).
\end{definition}

\begin{definition}[def:asymmetric-pure-stochastic-regime]\label{def:asymmetric-pure-stochastic-regime}
\par\smallskip\noindent\emph{Defined object.} \texttt{Asymmetric pure-\allowbreak{}stochastic regime}
\par\smallskip\noindent\emph{Construction.} \par\noindent \makebox[\linewidth][c]{\resizebox{0.98\linewidth}{!}{\(\displaystyle \mathcal R_n^{\rm asym}=\{(\bar\delta_{\mu,n},\bar\delta_{\pi,n}):\bar\delta_{\mu,n}=(0,s_{\mu,n}),\ \bar\delta_{\pi,n}=(0,s_{\pi,n}),\ n^{-1/2}\le s_{\pi,n}=o(s_{\mu,n}),\ s_{\mu,n}s_{\pi,n}\to0,\ n^{-1/2}=o(s_{\mu,n}s_{\pi,n})\}\)}}\par\noindent.
\par\smallskip\noindent\emph{Inputs.} \(n\); \(\bar\delta_{\mu,n}\); \(\bar\delta_{\pi,n}\); \(s_{\mu,n}\); \(s_{\pi,n}\).
\end{definition}

\begin{definition}[def:a1-displayed-benchmarks]\label{def:a1-displayed-benchmarks}
\par\smallskip\noindent\emph{Defined object.} \texttt{Theorem A.1 displayed benchmarks}
\par\smallskip\noindent\emph{Construction.} In Appendix A.1 of \cite{gu2026optimally}, the outcome error is denoted \(\varepsilon_Y\), and its condition is \(\mathbb E[\varepsilon_Y\mid X,T]=0\). Theorem A.1 displays \par\noindent \makebox[\linewidth][c]{\resizebox{0.98\linewidth}{!}{\(\displaystyle n^{-1/2}+[\delta_\mu^{\rm appr}+(\delta_\mu^{\rm stoc})^2]\wedge[(\delta_\mu^{\rm appr}+\delta_\mu^{\rm stoc})(\delta_\pi^{\rm appr}+\delta_\pi^{\rm stoc})]\)}}\par\noindent and \par\noindent \makebox[\linewidth][c]{\resizebox{0.98\linewidth}{!}{\(\displaystyle n^{-1/2}+[\delta_\mu^{\rm appr}+(\delta_\mu^{\rm stoc})^2][\delta_\pi^{\rm appr}+(\delta_\pi^{\rm stoc})^2]+(\delta_\mu^{\rm stoc}\wedge\delta_\pi^{\rm stoc})^2\)}}\par\noindent. At the proposed budgets these define \(\mathsf U_n=n^{-1/2}+\min\{s_{\mu,n}^2,s_{\mu,n}s_{\pi,n}\}\) and \(\mathsf L_n=n^{-1/2}+s_{\mu,n}^2s_{\pi,n}^2+s_{\pi,n}^2\). The paper uses \(2n\) observations, which changes only constants.
\par\smallskip\noindent\emph{Inputs.} \(n\); \(s_{\mu,n}\); \(s_{\pi,n}\); \(\mathsf U_n\); \(\mathsf L_n\).
\end{definition}

\section{Main results}

\begin{theorem}[thm:fixed-class-product-converse-fails-without-covariate-richness]\label{thm:fixed-class-product-converse-fails-without-covariate-richness}
Suppose \(\nu_d=\delta_{x_0}\) for some \(x_0\in[0,1]^d\). There is a numerical constant \(c_0>0\) such that, for every \(n\in\mathbb N\) and every pair of budgets \(\bar\delta_\mu,\bar\delta_\pi\in[0,1]^2\),
\[
\frac{c_0}{\sqrt{2n}}\le m_{\rm PL}^{0}(n,\bar\delta_\mu,\bar\delta_\pi)\le\frac{648}{\sqrt{2n}}.
\]
Moreover, for every sequence \((\bar\delta_{\mu,n},\bar\delta_{\pi,n})\in\mathcal R_n^{\rm asym}\) and every fixed \(F_{\mu,n}\in\mathcal H_n(0,s_{\mu,n})\) and \(F_{\pi,n}\in\mathcal H_n(0,s_{\pi,n})\),
\[
\makebox[\linewidth][c]{\resizebox{0.98\linewidth}{!}{\par\noindent \makebox[\linewidth][c]{\resizebox{0.98\linewidth}{!}{\(\displaystyle \displaystyle \inf_{\widehat\beta_n}\sup_{P\in\mathcal P_{\rm PL}^{0,\rm na}(F_{\mu,n},F_{\pi,n})}\int\left|\widehat\beta_n(o_{1:2n})-\beta(P)\right|\,dP^{\otimes2n}(o_{1:2n})\le\frac{648}{\sqrt{2n}}=o(s_{\mu,n}s_{\pi,n}).\)}}\par\noindent}}
\]
Consequently, the proposed fixed-class product converse is false under the frozen core: a product-scale converse requires an additional covariate-richness condition.
\end{theorem}
\begin{proof}
Put \(m=2n\). We first prove the upper bound uniformly over the full model. Fix arbitrary \(F_\mu,F_\pi\subseteq\mathcal T_3\) and \(P=P^0_{\beta,v,\mu,\pi}\in\mathcal P_{\rm PL}^{0}(F_\mu,F_\pi)\). Because \(X=x_0\) almost surely, \(\pi(X)=\pi(x_0)\) and \(\mu(X)=\mu(x_0)\) are constants. The treatment equation and its conditional moments imply
\[
\mathbb E_P[T]=\pi(x_0),\qquad V:=\operatorname{Var}_P(T)=\mathbb E_P[u^2]=v\ge\frac13.
\]
The outcome-error condition gives \(\mathbb E_P[\varepsilon\mid T]=0\), and therefore
\[
C:=\operatorname{Cov}_P(T,Y)=\operatorname{Cov}_P\bigl(T,\beta T+\mu(x_0)+\varepsilon\bigr)=\beta V.
\]
Thus the causal coefficient is the identified covariance ratio \(\beta(P)=C/V\).

For the observed sample define
\[
\overline T=m^{-1}\sum_{i=1}^mT_i,\qquad \overline Y=m^{-1}\sum_{i=1}^mY_i,
\]
\[
\makebox[\linewidth][c]{\resizebox{0.98\linewidth}{!}{\par\noindent \makebox[\linewidth][c]{\resizebox{0.98\linewidth}{!}{\(\displaystyle \displaystyle \widehat C=m^{-1}\sum_{i=1}^mT_iY_i-\overline T\,\overline Y,\qquad \widehat V=m^{-1}\sum_{i=1}^mT_i^2-\overline T^2,\)}}\par\noindent}}
\]
and the measurable estimator
\[
\makebox[\linewidth][c]{\resizebox{0.98\linewidth}{!}{\par\noindent \makebox[\linewidth][c]{\resizebox{0.98\linewidth}{!}{\(\displaystyle \displaystyle \widehat\beta^{\rm deg}=\Pi_{[-3,3]}\!\left(\frac{\widehat C}{\widehat V\vee(1/6)}\right),\qquad \Pi_{[-3,3]}(z)=\max\{-3,\min\{3,z\}\}.\)}}\par\noindent}}
\]
If \(W_1,\ldots,W_m\) are independent copies of a random variable with \(|W|\le B\), then Cauchy--Schwarz and independence give
\[
\makebox[\linewidth][c]{\resizebox{0.98\linewidth}{!}{\par\noindent \makebox[\linewidth][c]{\resizebox{0.98\linewidth}{!}{\(\displaystyle \displaystyle \mathbb E\left|m^{-1}\sum_{i=1}^mW_i-\mathbb EW\right|\le\sqrt{\frac{\operatorname{Var}(W)}m}\le\frac{B}{\sqrt m}.\)}}\par\noindent}}
\]
Apply this inequality to \(T\), \(Y\), \(TY\), and \(T^2\), whose absolute bounds are \(3,3,9,9\), respectively. Since \(|\overline T|\vee|\overline Y|\vee|\mathbb ET|\vee|\mathbb EY|\le3\),
\[
\makebox[\linewidth][c]{\resizebox{0.98\linewidth}{!}{\par\noindent \makebox[\linewidth][c]{\resizebox{0.98\linewidth}{!}{\(\displaystyle \displaystyle \mathbb E\left|\overline T\,\overline Y-(\mathbb ET)(\mathbb EY)\right|\le3\mathbb E|\overline Y-\mathbb EY|+3\mathbb E|\overline T-\mathbb ET|\le\frac{18}{\sqrt m}.\)}}\par\noindent}}
\]
Consequently,
\[
\mathbb E|\widehat C-C|\le\frac{27}{\sqrt m}.
\]
Likewise, \(|\overline T^2-(\mathbb ET)^2|\le6|\overline T-\mathbb ET|\), so
\[
\mathbb E|\widehat V-V|\le\frac{27}{\sqrt m}.
\]
Let \(D=\widehat V\vee(1/6)\). The map \(z\mapsto z\vee(1/6)\) is one-Lipschitz and fixes \(V\ge1/3\); hence \(|D-V|\le|\widehat V-V|\) and \(D\ge1/6\). Since \(C=\beta V\) and \(|\beta|\le3\),
\[
\makebox[\linewidth][c]{\resizebox{0.98\linewidth}{!}{\par\noindent \makebox[\linewidth][c]{\resizebox{0.98\linewidth}{!}{\(\displaystyle \displaystyle \left|\frac{\widehat C}{D}-\beta\right|\le\frac{|\widehat C-C|}{D}+\frac{|\beta|\,|V-D|}{D}\le6|\widehat C-C|+18|\widehat V-V|.\)}}\par\noindent}}
\]
Projection onto \([-3,3]\) cannot increase distance from \(\beta\in[-3,3]\). Taking expectations yields
\[
\makebox[\linewidth][c]{\resizebox{0.98\linewidth}{!}{\par\noindent \makebox[\linewidth][c]{\resizebox{0.98\linewidth}{!}{\(\displaystyle \displaystyle \mathbb E_{P^{\otimes m}}|\widehat\beta^{\rm deg}-\beta(P)|\le6\frac{27}{\sqrt m}+18\frac{27}{\sqrt m}=\frac{648}{\sqrt m}.\)}}\par\noindent}}
\]
This bound used no non-atomicity and no property of the fixed classes. The same estimator therefore bounds the inner minimax risk for every fixed class pair. Taking the two outer suprema in the definition of \(m_{\rm PL}^{0}\) proves
\[
m_{\rm PL}^{0}(n,\bar\delta_\mu,\bar\delta_\pi)\le\frac{648}{\sqrt m}.
\]

We next prove a matching lower bound inside an explicit bounded non-atomic submodel. Take
\[
F_\mu=F_\pi=G_\mu=G_\pi=\{0\}.
\]
The approximation error is zero. Also \(\partial G_\mu=\partial G_\pi=\{0\}\), whose localized Rademacher complexity is zero at every positive radius. Therefore
\[
\makebox[\linewidth][c]{\resizebox{0.98\linewidth}{!}{\par\noindent \makebox[\linewidth][c]{\resizebox{0.98\linewidth}{!}{\(\displaystyle \displaystyle \sup_{\delta>0:\,\delta\ge\delta^{\rm stoc}}\frac{\mathcal R_{n,\nu_d}(\delta;\{0\})}{3\delta}=0\le\delta^{\rm stoc}\)}}\par\noindent}}
\]
for every stochastic budget, including \(\delta^{\rm stoc}=0\). Thus the singleton truth classes belong to \(\mathcal H_n(\bar\delta_\mu)\) and \(\mathcal H_n(\bar\delta_\pi)\) for every budget pair.

Define
\[
b(z)=\exp\!\left\{-\frac{1}{1-z^2}\right\}\mathbf 1\{|z|<1\},\qquad Z=\int_{-1}^1b(z)^2\,dz,
\]
\[
q(z)=\frac{b(z)^2}{Z},\qquad r(z)=\sqrt{q(z)},\qquad J=\int_{\mathbb R}|r'(z)|^2\,dz.
\]
The density \(q\) is symmetric, mean zero, \(C^\infty\), and supported on \([-1,1]\). The standard bump vanishes to every order at \(\{-1,1\}\), so \(0<J<\infty\). Set
\[
\kappa=\min\left\{1,\frac{\sqrt3}{8\sqrt J}\right\},\qquad a_m=\frac{\kappa}{\sqrt m}.
\]
Let \(T\sim\operatorname{Unif}[-1,1]\), let \(\varepsilon\) have density \(q\) independently of \(T\), and, for \(\sigma\in\{-1,1\}\), define \(P_\sigma\) by
\[
X=x_0,\qquad \mu=\pi=0,\qquad u=T,\qquad \beta_\sigma=\sigma a_m,\qquad Y=\sigma a_mT+\varepsilon.
\]
Then \(v=\mathbb E[T^2]=1/3\), \(\mathbb E[\varepsilon\mid X,T]=0\), and \(T\) is non-atomic. Moreover, \(|T|\le1\) and \(|Y|\le1+a_m\le2<3\). Hence
\[
\makebox[\linewidth][c]{\resizebox{0.98\linewidth}{!}{\par\noindent \makebox[\linewidth][c]{\resizebox{0.98\linewidth}{!}{\(\displaystyle \displaystyle P_\sigma\in\mathcal P_{\rm PL}^{0,\rm na}(\{0\},\{0\})\subseteq\mathcal P_{\rm PL}^{0}(\{0\},\{0\}),\qquad \beta(P_\sigma)=\sigma a_m.\)}}\par\noindent}}
\]

For probability laws with densities \(p\) and \(q\), write
\[
H^2(p,q)=\int(\sqrt p-\sqrt q)^2.
\]
For any real \(h\), the fundamental theorem of calculus and Minkowski's integral inequality give the translation bound
\[
\|r(\cdot-h)-r\|_2\le|h|\,\|r'\|_2.
\]
Conditioning on the common uniform treatment distribution and applying this inequality with translation length \(2a_m|t|\), we obtain
\[
\makebox[\linewidth][c]{\resizebox{0.98\linewidth}{!}{\par\noindent \makebox[\linewidth][c]{\resizebox{0.98\linewidth}{!}{\(\displaystyle \displaystyle H^2(P_1,P_{-1})=\frac12\int_{-1}^1\|r(\cdot-a_mt)-r(\cdot+a_mt)\|_2^2\,dt\le4a_m^2J\,\mathbb E[T^2]=\frac{4a_m^2J}{3}.\)}}\par\noindent}}
\]
If \(\rho(P,Q)=\int\sqrt{dP\,dQ}\) denotes Hellinger affinity, then \(H^2(P,Q)=2-2\rho(P,Q)\) and \(\rho(P^{\otimes m},Q^{\otimes m})=\rho(P,Q)^m\). The inequality \(1-z^m\le m(1-z)\) for \(z\in[0,1]\) therefore yields
\[
H^2(P_1^{\otimes m},P_{-1}^{\otimes m})\le mH^2(P_1,P_{-1})\le\frac{4\kappa^2J}{3}\le\frac1{16}.
\]
Cauchy--Schwarz applied to \(|p-q|=|\sqrt p-\sqrt q|(\sqrt p+\sqrt q)\) shows that total variation is bounded by \(H\). Consequently,
\[
\operatorname{TV}(P_1^{\otimes m},P_{-1}^{\otimes m})\le\frac14.
\]
For any estimator \(\widehat\beta\), let \(\phi=\mathbf 1\{\widehat\beta\ge0\}\). On \(\{\phi=0\}\), the loss under \(P_1\) is at least \(a_m\); on \(\{\phi=1\}\), the loss under \(P_{-1}\) is at least \(a_m\). Hence
\[
\makebox[\linewidth][c]{\resizebox{0.98\linewidth}{!}{\par\noindent \makebox[\linewidth][c]{\resizebox{0.98\linewidth}{!}{\(\displaystyle \displaystyle \max_{\sigma\in\{-1,1\}}\mathbb E_{P_\sigma^{\otimes m}}|\widehat\beta-\sigma a_m|\ge\frac{a_m}{2}\left\{P_1^{\otimes m}(\phi=0)+P_{-1}^{\otimes m}(\phi=1)\right\}.\)}}\par\noindent}}
\]
The sum of the two testing errors is at least one minus total variation, so
\[
\makebox[\linewidth][c]{\resizebox{0.98\linewidth}{!}{\par\noindent \makebox[\linewidth][c]{\resizebox{0.98\linewidth}{!}{\(\displaystyle \displaystyle \max_{\sigma\in\{-1,1\}}\mathbb E_{P_\sigma^{\otimes m}}|\widehat\beta-\sigma a_m|\ge\frac{a_m}{2}\left(1-\frac14\right)=\frac{3\kappa}{8\sqrt m}.\)}}\par\noindent}}
\]
Taking the infimum over estimators and using the admissible singleton class pair in the two outer suprema proves the lower bound with \(c_0=3\kappa/8\).

Finally, fix a sequence in \(\mathcal R_n^{\rm asym}\). For every fixed admissible class pair, the full-model upper bound just proved also bounds the risk over its non-atomic subclass. The regime condition \(n^{-1/2}=o(s_{\mu,n}s_{\pi,n})\) gives
\[
\frac{648/\sqrt{2n}}{s_{\mu,n}s_{\pi,n}}\longrightarrow0.
\]
Thus every fixed-pair non-atomic risk is \(o(s_{\mu,n}s_{\pi,n})\), contradicting the proposed positive product lower bound for all sufficiently large \(n\). This proves the negative answer and identifies absent covariate richness as the obstruction.
\end{proof}

\section{Interpretation}

When \(X\) is degenerate, the partially linear model becomes a bounded one-regressor problem with treatment variance bounded away from zero. The nuisance budgets then do not affect the minimax order: the coefficient is a covariance ratio and its risk has matching root-\(n\) upper and lower bounds. Covariate variation is therefore essential to any slower product-scale frontier.

\section*{Technical internal limitation (diagnostic only)}

The verified arXiv:2606.06368v1 source, \texttt{append.tex}, line 18 uses \(\varepsilon_Y\), with \(Y\) as the outcome-error subscript, and imposes \(\mathbb E[\varepsilon_Y\mid X,T]=0\). The genuine limitation is that the matching minimax-order theorem treats the degenerate-covariate boundary and does not determine the minimax rate under a separately imposed covariate-richness condition.

\section{Honest scope}

The theorem proves the matching root-\(n\) minimax order only for \(\nu_d=\delta_{x_0}\). It shows that a slower product-scale lower bound cannot be uniform over the frozen design class, but it does not establish or refute such a lower bound within a quantitatively rich design subclass. It does not tighten Gu, Yin, Cai, and Fan's displayed lower bound.

\begin{thebibliography}{99}

  \bibitem{robinson1988root} Robinson, P. M. (1988). Root-\(N\)-consistent semiparametric regression. Econometrica 56, 931--954.

  \bibitem{chernozhukov2018dml} Chernozhukov, V., Chetverikov, D., Demirer, M., Duflo, E., Hansen, C., Newey, W., and Robins, J. (2018). Double/debiased machine learning for treatment and structural parameters. The Econometrics Journal 21, C1--C68.

  \bibitem{gu2025open} Gu, Y. (2025). Open Problem: Structure-Agnostic Minimax Risk for Partial Linear Model. Proceedings of Machine Learning Research 291, 6220--6224.

  \bibitem{gu2026optimally} Gu, Y., Yin, Q., Cai, T., and Fan, J. (2026). Optimally taming biases in black-box models for efficient semiparametric estimation. arXiv:2606.06368v1.

  \bibitem{jin2025hardnormal} Jin, J., Mackey, L., and Syrgkanis, V. (2025). It's Hard to Be Normal: The Impact of Noise on Structure-agnostic Estimation. Advances in Neural Information Processing Systems 38; arXiv:2507.02275.

  \bibitem{jin2025sharp} Jin, J. and Syrgkanis, V. (2025). Sharp Structure-Agnostic Lower Bounds for General Linear Functional Estimation. arXiv:2512.17341.

  \bibitem{balakrishnan2023fundamental} Balakrishnan, S., Kennedy, E. H., and Wasserman, L. (2023). The Fundamental Limits of Structure-Agnostic Functional Estimation. arXiv:2305.04116.

\end{thebibliography}

\end{document}